\documentclass[10pt,a4paper]{amsart}
\usepackage[margin=19mm]{geometry}
\usepackage{amsmath,amssymb,mathtools}
\usepackage[scr=rsfs,cal=euler]{mathalfa}
\usepackage{xfrac}
\usepackage[unicode,hidelinks,bookmarks=false]{hyperref}
\newcommand{\ie}{i.e.\ }
\newcommand{\cf}{cf.\ }
\newcommand{\resp}{respectively\ }
\newcommand{\Subsection}{\subsection}
\providecommand{\nobreakdash}{\nobreak\hbox{-}}
\setlength{\emergencystretch}{2em}
\multlinegap0pt
\usepackage{mathtools}
\usepackage{amssymb}
\usepackage{float}
\newtheorem{theorem}{Theorem}
\title{On the conservation of helicity by weak solutions of the 3D Euler and inviscid MHD equations}
\author{Daniel W. Boutros, Edriss S. Titi}
\date{Source: arXiv:2410.00813v2 (CC BY 4.0)}
\begin{document}
\maketitle

\noindent\emph{Source and licence.} On the conservation of helicity by weak solutions of the 3D Euler and inviscid MHD equations. Version arXiv:2410.00813v2 (CC BY 4.0), distributed by arXiv under the Creative Commons Attribution 4.0 International licence, http://creativecommons.org/licenses/by/4.0/, as recorded in the arXiv licence metadata for this exact version and shown on the abstract page https://arxiv.org/abs/2410.00813v2. Full licence notice: \texttt{licenses/arxiv-2410.00813-CC-BY-4.0.txt}.\par\smallskip

\noindent\textit{Editorial scope.} Counted unit: the article's theorem scalingthm (source0.tex lines 666--676). The article's complete original proof of it is delivered with it (source0.tex lines 678--734, one \texttt{proof} environment of 57 lines). No mathematics is written, repaired or re-derived: the statement and the proof are the article's own lines, in the article's own notation, and no retained line is substituted or re-flowed.  Retained uncounted as the source-local context the proof needs: lines 321--336; lines 385--389; lines 391--394; lines 587--589.  Nothing was omitted from any retained span, so the package carries no omission record.  No retained line contains a citation, so the wrapper carries no bibliography and no bibliography entry.  No retained line contains a figure environment, an image-inclusion command or drawing code, so no image is delivered and the package declares no asset dependency.  Editorial formatting only, in addition: an AMS \texttt{amsart} preamble that restates the article's own declarations and macros used by the retained lines, namely the environments \texttt{theorem} on separate counters, together with the standard packages those lines need.  The retained environments print in this document as Theorem 1 in order of appearance, whereas the article prints the same items with the numbering of its own full document.  The complete native source of the article as distributed in the arXiv e-print for this exact version is bundled unchanged as \texttt{sources/166109/source0.tex}.
\begin{theorem} \label{helicitybalance}
Let $(u,\omega)$ be a functional vorticity solution of the Euler equations such that $u \in L^3 ((0,T); W^{1/2+,3} (\mathbb{T}^3))$. Then the following equation of local helicity balance holds for all test functions $\varphi \in \mathcal{D} (\mathbb{T}^3 \times (0,T); \mathbb{R})$
\begin{align}
&\int_0^T \bigg[ - \bigg\langle 2 \omega u, \partial_t \varphi \bigg\rangle + \bigg\langle  2 \nabla \cdot (p \omega) + \nabla \cdot \big[ 2 (u \cdot \omega) u - \lvert u \rvert^2 \omega \big], \varphi \bigg\rangle \bigg] dt \nonumber \\
&+ \bigg\langle D_{1} (u, \omega) - \frac{1}{2} D_{2} (u, \omega) ,  \varphi \bigg\rangle_{x,t} = 0. \label{localhelicitybalance}
\end{align}
where the defect terms are given by
\begin{align}
D_1 (u,\omega) - \frac{1}{2} D_2 (u, \omega) &\coloneqq \lim_{\epsilon \rightarrow 0} \bigg[\int_{\mathbb{R}^3} \nabla \phi_\epsilon (\xi) \cdot \delta u (\xi ; x,t) (\delta \omega (\xi ;x,t) \cdot \delta u (\xi ; x, t) ) d \xi \nonumber \\
&- \frac{1}{2} \int_{\mathbb{R}^3} \nabla \phi_\epsilon (\xi) \cdot \delta \omega (\xi ; x,t) \lvert \delta u (\xi ;x,t) \rvert^2 d \xi \bigg],
\end{align}
where the convergence is in the sense of distributions and independent of the choice of the mollifier. Specifically, we have the following evolution equation for the helicity density
\begin{equation}
2 \partial_t (u \omega) + 2 \nabla \cdot (p \omega) + \nabla \cdot \big[2 (u \cdot \omega) u - \lvert u \rvert^2 \omega \big] + \bigg[ D_1 (u,\omega) - \frac{1}{2} D_2 (u, \omega) \bigg] = 0, \quad \text{in } \mathcal{D}' (\mathbb{T}^3 \times (0,T)).
\end{equation}
\end{theorem}

\begin{align}
&\int_0^T \bigg\langle \partial_t \big( \omega^\epsilon u + \omega u^\epsilon \big) + \nabla \cdot (p \omega^\epsilon + p^\epsilon \omega) + \nabla \cdot \big[ (u \cdot \omega^\epsilon) u + (\omega \cdot u^\epsilon) u - (u \cdot u^\epsilon) \omega \big], \varphi \bigg\rangle dt \nonumber \\
&+ \frac{1}{2} \int_0^T \bigg\langle \nabla \cdot \big[ 2((\omega \cdot u) u)^\epsilon - 2(\omega \cdot u)^\epsilon u - ( \lvert u \rvert^2 \omega)^\epsilon + \omega (\lvert u \rvert^2)^\epsilon \big], \varphi \bigg\rangle dt \nonumber \\
&= \int_0^T \bigg\langle - D_{1,\epsilon} (u, \omega) + \frac{1}{2} D_{2,\epsilon} (u, \omega) ,  \varphi \bigg\rangle dt. \label{mollifiedbalance}
\end{align}

\begin{align}
&\int_0^T \bigg[ - \bigg\langle 2 \omega u, \partial_t \varphi \bigg\rangle + \bigg\langle  2 \nabla \cdot (p \omega) + \nabla \cdot \big[ 2 (u \cdot \omega) u - \lvert u \rvert^2 \omega \big], \varphi \bigg\rangle \bigg] dt \nonumber \\
&= \bigg\langle - D_{1} (u, \omega) + \frac{1}{2} D_{2} (u, \omega) ,  \varphi \bigg\rangle_{x,t}. \label{limithelicitybalance}
\end{align}

\begin{equation} \label{scalinglaw}
\big\langle \delta u_L (\xi;\cdot,\cdot) \big( \delta u (\xi;\cdot,\cdot) \cdot \delta \omega (\xi;\cdot,\cdot) \big) \rangle - \frac{1}{2} \langle \delta \omega_L (\xi;\cdot,\cdot) \lvert \delta u(\xi;\cdot,\cdot) \rvert^2 \rangle = - \frac{4}{3} \tilde{\epsilon} \xi,
\end{equation}

\begin{theorem} \label{scalingthm}
Let $(u,\omega)$ be a functional vorticity solution of the Euler equations such that $u \in L^3 ((0,T);B^{1/2+}_{3,\infty} (\mathbb{T}^3))$. Then the following distributional limit converges in $W^{-1,1} ((0,T); \linebreak W^{-2,1} (\mathbb{T}^3))$
\begin{align}
&\lim_{r \rightarrow 0} \bigg[ \frac{S_1 (u,\omega,r)}{r} - \frac{1}{2} \frac{S_2 (u,\omega,r)}{r} \bigg].
\end{align}
Moreover, the following distributional equation holds in $W^{-1,1} ((0,T); W^{-2,1} (\mathbb{T}^3))$ (for any $\varphi \in \mathcal{D} (\mathbb{T}^3 \times (0,T);\mathbb{R})$)
\begin{equation}
-\frac{4}{3} \cdot \frac{1}{4} \bigg\langle D_{1} (u, \omega) - \frac{1}{2} D_{2} (u, \omega) , \varphi \bigg\rangle_{x,t} = \bigg\langle \lim_{r \rightarrow 0} \bigg[ \frac{1}{4 \pi} \frac{S_1 (u,\omega,r)}{r} - \frac{1}{2} \cdot \frac{1}{4 \pi} \frac{S_2 (u,\omega,r)}{r} \bigg] , \varphi \bigg\rangle_{x,t},
\end{equation}
which is the inviscid version of the scaling law given in equation \eqref{scalinglaw}.
\end{theorem}

\begin{proof}
Let us first define a new (nonincreasing) $C^\infty$ function $\kappa: \mathbb{R} \rightarrow \mathbb{R}$ by
\begin{equation*}
\kappa (s) = \begin{cases}
1 \quad \text{if } s \leq \frac{3}{4}, \\
0 \quad \text{if } s \geq \frac{5}{4}.
\end{cases}
\end{equation*}
We then define a $C^\infty$ function $\kappa_{\mu,\epsilon}: \mathbb{R}^3 \rightarrow \mathbb{R}$ (for $\mu, \epsilon > 0$)
\begin{equation*}
\widetilde{\kappa}_{\mu,\epsilon} (x) \coloneqq \frac{1}{\epsilon^3 } \kappa \left( 1 + \mu^{-1} \left(\frac{\lvert x \rvert}{\epsilon} -1 \right) \right).
\end{equation*}
Then we define the following smooth mollifier $\kappa_{\mu,\epsilon}$ as follows
\begin{equation*}
\kappa_{\mu,\epsilon} (x) \coloneqq \frac{1}{\lVert \widetilde{\kappa}_{\mu,1} \rVert_{L^1}} \widetilde{\kappa}_{\mu,\epsilon} (x).
\end{equation*}
Similarly to before, we will use the notation $u^{\mu,\epsilon} \coloneqq u * \kappa_{\mu,\epsilon}$. Next we observe that the derivation of equation \eqref{mollifiedbalance} (in the proof of the equation of local helicity balance in Theorem \ref{helicitybalance}) does not rely on any particular choice of mollifier, therefore equation \eqref{mollifiedbalance} also holds with $\phi_\epsilon$ replaced by $\kappa_{\mu,\epsilon}$, which leads to (for any $\varphi \in \mathcal{D} (\mathbb{T}^3 \times (0,T);\mathbb{R})$)
\begin{align}
&\int_0^T \bigg\langle \partial_t \big( \omega^{\mu,\epsilon} u + \omega u^{\mu,\epsilon} \big) + \nabla \cdot (p \omega^{\mu,\epsilon} + p^{\mu,\epsilon} \omega) + \nabla \cdot \big[ (u \cdot \omega^{\mu,\epsilon}) u + (\omega \cdot u^{\mu,\epsilon}) u - (u \cdot u^{\mu,\epsilon}) \omega \big], \varphi \bigg\rangle dt \nonumber \\
&+ \frac{1}{2} \int_0^T \bigg\langle \nabla \cdot \big[ 2((\omega \cdot u) u)^{\mu,\epsilon} - 2(\omega \cdot u)^{\mu,\epsilon} u - ( \lvert u \rvert^2 \omega)^{\mu,\epsilon} + \omega (\lvert u \rvert^2)^{\mu,\epsilon} \big], \varphi \bigg\rangle dt \nonumber \\
&= \int_0^T \bigg\langle - D_{1,\mu,\epsilon} (u, \omega) + \frac{1}{2} D_{2,\mu,\epsilon} (u, \omega) ,  \varphi \bigg\rangle dt, \label{mollifiedsphericalbalance}
\end{align}
where the defect terms are now given by
\begin{align*}
\bigg\langle D_{1,\mu,\epsilon} (u, \omega) - \frac{1}{2} D_{2,\mu,\epsilon} (u, \omega), \varphi \bigg\rangle &= \int_{\mathbb{R}^3} \nabla \kappa_{\mu,\epsilon} (\xi) \cdot \big\langle \delta u (\xi ; \cdot, t) (\delta \omega (\xi ; \cdot, t) \cdot \delta u (\xi ; \cdot, t) ), \varphi \big\rangle d \xi \\
&- \frac{1}{2} \int_{\mathbb{R}^3} \nabla \kappa_{\mu,\epsilon} (\xi) \cdot \big\langle \delta \omega (\xi ; \cdot, t) \lvert \delta u (\xi ; \cdot, t) \rvert^2, \varphi \big\rangle d \xi.
\end{align*}
Note that the defect terms $D_{1,\mu,\epsilon} (u, \omega)$ and $D_{2,\mu,\epsilon} (u, \omega)$ are elements of $L^1 ((0,T);B^{-1/2}_{1,\infty} (\mathbb{T}^3))$. Now because the mollifier $\kappa_{\mu,\epsilon}$ is radial, we will change to spherical coordinates and write (where we make the change of variable $r \xi' \coloneqq \frac{\xi}{\epsilon}$ where we choose $\lvert \xi' \rvert = 1$ and we will relabel $\xi'$ by $\xi$)
\begin{align*}
&\bigg\langle D_{1,\mu,\epsilon} (u, \omega) - \frac{1}{2} D_{2,\mu,\epsilon} (u, \omega), \varphi \bigg\rangle \\
&= \int_{0}^\infty dr \frac{1}{\lVert \widetilde{\kappa}_{\mu,1} \rVert_{L^1}} \kappa' ( 1 + \mu^{-1} (r -1) ) \frac{r^3}{\epsilon \mu r} \bigg[ \int_{\mathbb{S}^2} \big\langle \xi \cdot \delta u (\epsilon r \xi ; \cdot,t) \big(\delta \omega (\epsilon r \xi ;\cdot,t) \cdot \delta u (\epsilon r \xi ; \cdot, t) \big), \varphi \big\rangle d S (\xi) \\
&- \frac{1}{2} \int_{\mathbb{S}^2} \big\langle \xi \cdot \delta \omega (\epsilon r \xi ; \cdot,t) \lvert \delta u (\epsilon r \xi ;\cdot,t) \rvert^2, \varphi \big\rangle d S (\xi) \bigg] \\
&= \int_0^\infty dr \frac{1}{\lVert \widetilde{\kappa}_{\mu,1} \rVert_{L^1}} \kappa' ( 1 + \mu^{-1} (r -1) ) \frac{r^3}{\epsilon \mu} \bigg[ \frac{\langle S_1 (u,\omega, \epsilon r), \varphi \rangle - \frac{1}{2} \langle S_2 (u,\omega, \epsilon r), \varphi \rangle}{r} \bigg].
\end{align*}
We now show that the right-hand side converges to $- \frac{3}{4 \pi \epsilon} \langle S_1 (u,\omega, \epsilon), \varphi \rangle + \frac{3}{8 \pi \epsilon} \langle S_2 (u,\omega, \epsilon), \varphi \rangle$, subtracting these terms from the equation gives
\begin{align*}
&\bigg\lvert \int_0^\infty dr \frac{1}{\lVert \widetilde{\kappa}_{\mu,1} \rVert_{L^1}} \kappa' ( 1 + \mu^{-1} (r -1) ) \frac{r^3}{\epsilon \mu} \bigg[ \frac{\langle S_1 (u,\omega, \epsilon r), \varphi \rangle - \frac{1}{2} \langle S_2 (u,\omega, \epsilon r), \varphi \rangle}{r} - \langle S_1 (u,\omega, \epsilon ), \varphi \rangle \\
&+ \frac{1}{2} \langle S_2 (u,\omega, \epsilon ), \varphi \rangle \bigg] \bigg\rvert \\
&\lesssim \frac{1}{\epsilon \mu} \int_{1 - \frac{1}{4} \mu}^{1 + \frac{1}{4} \mu} r^3 \bigg\lvert \frac{\langle S_1 (u,\omega, \epsilon r), \varphi \rangle}{r} - \langle S_1 (u,\omega, \epsilon), \varphi \rangle - \frac{1}{2} \frac{\langle S_2 (u,\omega, \epsilon r), \varphi \rangle}{r} \\
&+ \frac{1}{2} \langle S_2 (u,\omega, \epsilon), \varphi \rangle \bigg\rvert dr \xrightarrow[]{\mu \rightarrow 0} 0,
\end{align*}
where we have used the continuity of the translation operator in Besov spaces. We observe that the left-hand side of equation \eqref{mollifiedsphericalbalance} converges in the sense of distributions as $\mu \rightarrow 0$. This is because $\lim_{\mu \rightarrow 0} \kappa_{\mu,\epsilon} = \epsilon^{-3} 1_{B_\epsilon (0)} \eqqcolon \kappa_{0,\epsilon}$ in $L^p (\mathbb{R}^3)$ for any $1 \leq p < \infty$ (for fixed $\epsilon > 0$), where $B_\epsilon (0)$ is the ball of radius $\epsilon$ around $0$. Therefore we find that
\begin{align}
&\int_0^T \bigg\langle \partial_t \big( \omega^{0,\epsilon} u + \omega u^{0,\epsilon} \big) + \nabla \cdot (p \omega^{0,\epsilon} + p^{0,\epsilon} \omega) + \nabla \cdot \big[ (u \cdot \omega^{0,\epsilon}) u + (\omega \cdot u^{0,\epsilon}) u - (u \cdot u^{0,\epsilon}) \omega \big], \varphi \bigg\rangle dt \nonumber \\
&+ \frac{1}{2} \int_0^T \bigg\langle \nabla \cdot \big[ 2((\omega \cdot u) u)^{0,\epsilon} - 2(\omega \cdot u)^{0,\epsilon} u - ( \lvert u \rvert^2 \omega)^{0,\epsilon} + \omega (\lvert u \rvert^2)^{0,\epsilon} \big], \varphi \bigg\rangle dt \nonumber \\
&= \frac{3}{4 \pi} \bigg\langle \frac{S_1 (u,\omega,\epsilon)}{\epsilon} - \frac{1}{2} \frac{S_2 (u,\omega,\epsilon)}{\epsilon} , \varphi \bigg\rangle_{x,t}. \label{sphericalaverageexpression}
\end{align}
We observe that as $\epsilon \rightarrow 0$ the left-hand side of equation \eqref{sphericalaverageexpression} converges in $W^{-1,1} ((0,T); W^{-2,1} (\mathbb{T}^3))$ to the left-hand side of equation \eqref{limithelicitybalance}. Therefore we can conclude that the limit of the spherical averages $\lim_{\epsilon \rightarrow 0} \bigg[ \frac{1}{\epsilon} S_1 (u,\omega,\epsilon) - \frac{1}{2 \epsilon} S_2 (u,\omega,\epsilon) \bigg]$ is well-defined as a distribution and is equal to the right-hand side of equation \eqref{limithelicitybalance}, which leads to the following equality
\begin{equation*}
\bigg\langle - D_{1} (u, \omega) + \frac{1}{2} D_{2} (u, \omega) ,  \varphi \bigg\rangle_{x,t} = \frac{3}{4 \pi} \bigg\langle \lim_{\epsilon \rightarrow 0} \bigg[ \frac{S_1 (u,\omega,\epsilon)}{\epsilon} - \frac{1}{2} \frac{S_2 (u,\omega,\epsilon)}{\epsilon} \bigg] , \varphi \bigg\rangle_{x,t}.
\end{equation*}
We observe that the left-hand side can be formally identified (up to a minus sign) with the average rate of change of the helicity in the vanishing viscosity limit if one divides by a factor of $4$. Moreover, we note that the factors $(4 \pi)^{-1}$ are part of the spherical averages. Therefore we can rewrite this relation as follows
\begin{equation}
-\frac{4}{3} \cdot \frac{1}{4} \bigg\langle D_1 (u, \omega) - \frac{1}{2} D_2 (u, \omega), \varphi \bigg\rangle_{x,t} = \bigg\langle \lim_{\epsilon \rightarrow 0} \bigg[ \frac{1}{4 \pi} \frac{S_1 (u,\omega,\epsilon)}{\epsilon} - \frac{1}{2} \cdot \frac{1}{4 \pi} \frac{S_2 (u,\omega,\epsilon)}{\epsilon} \bigg] , \varphi \bigg\rangle_{x,t},
\end{equation}
which holds in $W^{-1,1} ((0,T); W^{-2,1} (\mathbb{T}^3))$ and is in agreement with equation \eqref{scalinglaw}, and therefore concludes the proof.
\end{proof}

\end{document}
